\documentclass[10pt]{article}
\usepackage[a4paper,margin=0.86in]{geometry}
\renewcommand{\rmdefault}{cmr}
\usepackage{amsmath,amssymb,amsthm,mathtools,amscd,enumitem,xcolor,hyperref}
\hypersetup{colorlinks=true,linkcolor=blue!55!black,citecolor=blue!55!black,urlcolor=blue!55!black}
\setlength{\emergencystretch}{4em}
\setlength{\parindent}{0pt}
\setlength{\parskip}{0.52em}
\setlist{leftmargin=1.55em,itemsep=0.2em,topsep=0.25em}
\sloppy

\newtheorem{theorem}{Theorem}[section]
\newtheorem{proposition}[theorem]{Proposition}
\newtheorem{lemma}[theorem]{Lemma}
\newtheorem{corollary}[theorem]{Corollary}
\theoremstyle{definition}
\newtheorem{definition}[theorem]{Definition}
\newtheorem{convention}[theorem]{Convention}
\newtheorem{remark}[theorem]{Remark}

\newcommand{\Zl}{\mathbf Z_\ell}
\newcommand{\Ql}{\mathbf Q_\ell}
\newcommand{\Qlb}{\overline{\mathbf Q}_\ell}
\newcommand{\OO}{\mathcal O}
\newcommand{\Dbc}{D^b_c}
\newcommand{\Dctf}{D^b_{\mathrm{ctf}}}
\newcommand{\HomD}{\operatorname{Hom}}
\newcommand{\red}{\operatorname{red}}

\begin{document}
\begin{center}
{\large\bfseries The Exceptional Right Adjoint \(Rf^!\)\par}
\medskip
{\normalsize Brown representability, compactifications, and coefficient descent\par}
\end{center}

\begin{abstract}
Let \(f:X\to Y\) be a separated morphism of finite type.  This note specifies
the ordinary triangulated categories on which \(Rf_!\) and its exceptional
right adjoint are defined.  It separates the representability input from the
construction of the adjunction, records the compactification formula and its
mates, and gives the comparison theorems needed for composition and change of
coefficients.  Every restricted functor is preceded by a preservation
statement.
\end{abstract}

\tableofcontents

\section{Coefficient stages and derived categories}

Fix a prime \(\ell\).  All schemes are noetherian of finite Krull dimension,
\(\ell\) is invertible on them, and all morphisms are quasi-compact and
quasi-separated.  The topology is the small pro-etale topology.  For finite
coefficients it may be replaced by the small etale topology through the
finite-coefficient comparison theorem.

\begin{definition}[Finite coefficient stages]
Let \(\Lambda\) be a finite local \(\mathbf Z/\ell^n\mathbf Z\)-algebra.  Write

\[
 D(X,\Lambda)=D\bigl(\mathsf{Mod}(X_{\mathrm{pro\acute et}},\Lambda_X)\bigr)
\]

for the Verdier localization of the category of cochain complexes of sheaves
of \(\Lambda_X\)-modules at quasi-isomorphisms.  Define
\(\Dbc(X,\Lambda)\) to be the full triangulated subcategory whose objects have
only finitely many nonzero cohomology sheaves and whose cohomology sheaves are
locally constant with finite stalks on the members of a finite constructible
stratification.  Define \(\Dctf(X,\Lambda)\) by imposing in addition finite
Tor-amplitude over \(\Lambda_X\).
For integers \(a\leq b\), write \(D^{[a,b]}_c(X,\Lambda)\) for the full
subcategory of \(\Dbc(X,\Lambda)\) whose cohomology sheaves vanish in
degrees outside \([a,b]\).
\end{definition}

\begin{definition}[Adic and field stages]
Let \(E/\Ql\) be finite, let \(\OO_E\) be its ring of integers, choose a
uniformizer \(\varpi\), and put \(\Lambda_n=\OO_E/\varpi^n\).
Let \(\mathsf{Mod}(X,\Lambda_\bullet)\) be the abelian category of inverse
systems \((M_n)\) of sheaves, where \(M_n\) is a
\(\Lambda_{n,X}\)-module and the transition maps are linear through
\(\Lambda_{n+1}\to\Lambda_n\).  Put

\[
 D(X,\Lambda_\bullet)=D\bigl(\mathsf{Mod}(X,\Lambda_\bullet)\bigr).
\]

Evaluation at level \(n\) is an exact functor on the system abelian
categories and hence an exact triangulated functor

\[
 \operatorname{ev}_{n,X}:D(X,\Lambda_\bullet)\longrightarrow
 D(X,\Lambda_n).
\]

For \(K\in D(X,\Lambda_\bullet)\), write
\(K_n=\operatorname{ev}_{n,X}K\).  The transition maps of the underlying
system induce, for \(m\leq n\), canonical normalization morphisms

\[
 \nu^K_{n,m}:
 K_n\otimes^L_{\Lambda_{n,X}}\Lambda_{m,X}\longrightarrow K_m.
\]

The object \(K\) is normalized if every \(\nu^K_{n,m}\) is an isomorphism.
These morphisms satisfy \(\nu^K_{n,n}=\operatorname{id}\) and

\[
 \nu^K_{n,r}=\nu^K_{m,r}\circ
 \bigl(\nu^K_{n,m}\otimes^L_{\Lambda_{m,X}}\Lambda_{r,X}\bigr)
 \qquad(r\leq m\leq n),
\]

under the canonical reassociation of derived tensor products.  Morphisms of
normalized systems are morphisms in \(D(X,\Lambda_\bullet)\), represented by
roofs of maps of system complexes; they are not arbitrary families of
levelwise maps.  The category
\(\Dbc(X,\OO_E)\) is the full triangulated subcategory of normalized systems
for which all
\(K_n\) lie in \(\Dbc(X,\Lambda_n)\), with one common cohomological interval
and compatible finite stratifications.  The category
\(\Dctf(X,\OO_E)\) also requires a common finite Tor-amplitude interval.
The rational categories are the idempotent-complete essential images

\[
 \Dbc(X,E)=\operatorname{Kar}\bigl(\Dbc(X,\OO_E)[1/\varpi]\bigr),
 \qquad
 \Dctf(X,E)=\operatorname{Kar}\bigl(\Dctf(X,\OO_E)[1/\varpi]\bigr).
\]

Finally, \(\Dbc(X,\Qlb)\) and \(\Dctf(X,\Qlb)\) are the Karoubi completions of
the filtered \(2\)-colimits of these categories over finite extensions
\(E/\Ql\).  Thus every object and every morphism over \(\Qlb\) is represented
over some finite coefficient field, and every idempotent is retained as part
of the object.
\end{definition}

\begin{definition}[Reduction and scalar extension]
For \(m\leq n\), derived extension of scalars gives

\[
 \red_{n,m,X}:D(X,\Lambda_n)\longrightarrow D(X,\Lambda_m),
 \qquad K\longmapsto K\otimes^L_{\Lambda_{n,X}}\Lambda_{m,X}.
\]

Restriction of scalars along \(\Lambda_n\to\Lambda_m\) is the exact functor

\[
 \operatorname{res}_{m,n,X}:D(X,\Lambda_m)\longrightarrow D(X,\Lambda_n).
\]

On the ambient derived categories one has
\(\red_{n,m,X}\dashv\operatorname{res}_{m,n,X}\), with unit and counit
induced by the coefficient homomorphism.  No preservation of finite
Tor-amplitude by \(\operatorname{res}_{m,n,X}\) is asserted.

At the integral stage write

\[
 \red_{m,X}=\operatorname{ev}_{m,X}:
 D(X,\Lambda_\bullet)\longrightarrow D(X,\Lambda_m).
\]

It restricts to the displayed bounded constructible and finite-Tor
subcategories.

For a homomorphism \(A\to B\) between finite or finite-extension coefficient
rings, put

\[
 L\operatorname{Ext}_{B/A,X}(K)=K\otimes^L_{A_X}B_X.
\]

This is an exact triangulated functor on the finite-Tor subcategories.
\end{definition}

\section{The already constructed functor \texorpdfstring{\(Rf_!\)}{Rf!}}

\begin{definition}[Compactifiable morphism]
A morphism \(f:X\to Y\) is compactifiable if it admits a factorization

\[
 c=(X\xrightarrow{j_c}\overline X_c\xrightarrow{p_c}Y)
\]

with \(j_c\) an open immersion and \(p_c\) proper.  A refinement
\(q:c'\to c\) is a proper \(Y\)-morphism
\(q:\overline X_{c'}\to\overline X_c\) satisfying
\(qj_{c'}=j_c\).
\end{definition}

\begin{definition}[Extension by zero]
For a quasi-compact open immersion \(j:U\hookrightarrow X\), extension by
zero is the sheaf functor

\[
 j_!: \mathsf{Mod}(U,\Lambda)\longrightarrow
       \mathsf{Mod}(X,\Lambda)
\]

characterized by \(j^*j_!\cong\operatorname{id}\) and by zero restriction
to \(X\setminus U\).  It is exact and is left adjoint to
\(j^*:D(X,\Lambda)\to D(U,\Lambda)\).  Its degreewise application therefore
defines an exact triangulated functor, still denoted

\[
 j_!:D(U,\Lambda)\longrightarrow D(X,\Lambda).
\]
\end{definition}

\begin{proposition}[Open-extension finiteness]
The functor \(j_!\) preserves arbitrary coproducts, boundedness,
constructibility, and finite Tor-amplitude.  More precisely,

\[
 K\in D^{[a,b]}_c(U,\Lambda)
 \quad\Longrightarrow\quad
 j_!K\in D^{[a,b]}_c(X,\Lambda),
\]

and \(\mathcal H^r(j_!K)\cong j_!\mathcal H^r(K)\).
The same statements hold levelwise for normalized adic systems.
\end{proposition}

\begin{definition}[Proper derived direct image]
For a proper morphism \(p:Z\to Y\), choose a functorial K-injective
replacement \(K\to I(K)\) in the abelian category of sheaf complexes.  Set

\[
 Rp_*K=p_*I(K),
 \qquad Rp_*:D(Z,\Lambda)\longrightarrow D(Y,\Lambda).
\]

This definition is independent of the replacement up to the unique natural
isomorphism compatible with derived localization.
\end{definition}

\begin{theorem}[Proper finiteness and amplitude]
Let \(p:Z\to Y\) be proper of relative dimension at most \(d\).  Then

\[
 Rp_*:D^{[a,b]}_c(Z,\Lambda)\longrightarrow
 D^{[a,b+2d]}_c(Y,\Lambda).
\]

It restricts to exact triangulated functors on \(\Dbc\) and \(\Dctf\).
The formation of \(Rp_*\) commutes with finite reduction and with arbitrary
base change on bounded constructible objects.
\end{theorem}

\begin{definition}[Compactification candidate]
For a compactification \(c=(j_c,p_c)\), define

\[
 Rf_{!,c}=Rp_{c,*}j_{c,!}:D(X,\Lambda)\longrightarrow D(Y,\Lambda).
\]

For a refinement \(q:c'\to c\), the proper equality
\(p_{c'}=p_cq\), the identification
\(Rq_*j_{c',!}\cong j_{c,!}\), and the composition isomorphism for proper
direct image define

\[
 \rho_{q}:Rf_{!,c'}\xrightarrow{\sim}Rf_{!,c}.
\]
\end{definition}

\begin{theorem}[Compactification coherence for \(Rf_!\)]
The refinement maps satisfy

\[
 \rho_{\operatorname{id}_c}=\operatorname{id},
 \qquad \rho_{q q'}=\rho_q\circ\rho_{q'}.
\]

Any two compactifications have a common refinement.  Consequently there is
an exact triangulated functor

\[
 Rf_!:D(X,\Lambda)\longrightarrow D(Y,\Lambda)
\]

equipped with canonical isomorphisms \(Rf_!\cong Rf_{!,c}\), compatible with
all refinements.  It restricts to

\[
 Rf_!:D^{[a,b]}_c(X,\Lambda)\longrightarrow
 D^{[a,b+2d]}_c(Y,\Lambda)
\]

whenever a compactification has proper relative dimension at most \(d\).
For composable compactifiable maps there is a natural isomorphism

\[
 \alpha_{g,f}:Rg_!Rf_!\xrightarrow{\sim}R(gf)_!.
\]
For three composable maps the two composites of the \(\alpha\)'s are equal,
and for an identity map both unit composites are the identity.  These are
equalities of natural transformations between ordinary triangulated
functors.
\end{theorem}

\section{Representability and the finite exceptional inverse image}

\begin{definition}[Exceptional Hom functor]
Fix \(L\in D(Y,\Lambda)\).  Define

\[
 H_{f,L}:D(X,\Lambda)^{\mathrm{op}}\longrightarrow\mathsf{Ab},
 \qquad
 H_{f,L}(K)=\HomD_{D(Y,\Lambda)}(Rf_!K,L).
\]

On a morphism \(u:K\to K'\), the map \(H_{f,L}(u)\) is precomposition by
\(Rf_!(u)\).
\end{definition}

\begin{proposition}[Cohomologicality of the exceptional Hom functor]
For \(n\in\mathbf Z\), set

\[
 H^n_{f,L}(K)=
 \operatorname{Hom}_{D(Y,\Lambda)}(Rf_!K,L[n]).
\]

For every distinguished triangle
\(K'\to K\to K''\to K'[1]\), these groups form the contravariant long
exact sequence

\[
 \cdots\longrightarrow H^n_{f,L}(K'')
 \longrightarrow H^n_{f,L}(K)
 \longrightarrow H^n_{f,L}(K')
 \longrightarrow H^{n+1}_{f,L}(K'')
 \longrightarrow\cdots .
\]
\end{proposition}

\begin{definition}[Weakly contractible objects and generator sheaves]
An object \(U\to X\) of the small pro-etale site is \emph{weakly
contractible} if every epimorphism of sheaves of sets \(F\twoheadrightarrow G\)
induces a surjection \(F(U)\twoheadrightarrow G(U)\).  If \(U\) is also
affine and quasi-compact, let

\[
 G_{U,r}=u_!\Lambda_U[r]\in D(X,\Lambda),\qquad r\in\mathbf Z,
\]

where \(u:U\to X\), and where \(u_!\Lambda_U\) is the sheafification of the
free \(\Lambda\)-module presheaf on the representable presheaf \(h_U\).
Thus, for every sheaf \(F\),

\[
 \operatorname{Hom}_{\Lambda_X}(u_!\Lambda_U,F)\cong F(U).
\]
\end{definition}

\begin{theorem}[Small weakly contractible covering theorem]
Every object of the selected small pro-etale site of \(X\) has a small
covering family by affine, quasi-compact, weakly contractible objects.  The
family of all chosen objects may be taken to be a set, and two chosen
covering families have a common chosen refinement.
\end{theorem}

\begin{proposition}[Exactness and sums of sections]
If \(U\) is affine, quasi-compact, and weakly contractible, then
\(F\mapsto F(U)\) is exact and commutes with set-indexed direct sums.  Hence
for every \(K\in D(X,\Lambda)\),

\[
 \HomD_{D(X,\Lambda)}(G_{U,r},K)\cong H^{-r}(K(U)),
\]

and each \(G_{U,r}\) is compact.
\end{proposition}

\begin{theorem}[Compact generation of the ambient category]
The set
\(
 \mathcal G_X=\{G_{U,r}:U\text{ chosen as above},\ r\in\mathbf Z\}
\)
is stable under shifts and compactly generates \(D(X,\Lambda)\).  Explicitly,

\[
 \HomD(G_{U,r},\textstyle\bigoplus_iK_i)
 \cong\bigoplus_i\HomD(G_{U,r},K_i),
\]

and \(K=0\) if \(\HomD(G,K)=0\) for every \(G\in\mathcal G_X\).
\end{theorem}

\begin{definition}[Uniform lower bound and full cohomological dimension]
\label{def:full-cd}
A family \((K_i)_{i\in I}\) in \(D(Z,\Lambda)\) is uniformly bounded below
if there is \(a\in\mathbf Z\) such that \(H^q(K_i)=0\) for all \(i\) and
all \(q<a\).

Choose a set \(\mathcal B_Y\) of affine, quasi-compact, weakly contractible
objects covering the selected small pro-etale site of \(Y\).  For a proper
map \(p:Z\to Y\), write \(\operatorname{FullCD}(p,\Lambda;c)\) if, for every
sheaf \(F\) of \(\Lambda\)-modules on \(Z_{\mathrm{pro\acute et}}\),

\[
 R^qp_*F=0\qquad(q>c),
\]

and the same integer works after every selected site-local base change.
Explicitly, for \(V\in\mathcal B_Y\), put
\(p_V:Z\times_YV\to V\); the datum also requires

\[
 R^q(p_V)_*G=0\qquad(q>c)
\]

for every sheaf \(G\) of \(\Lambda\)-modules on
\((Z\times_YV)_{\mathrm{pro\acute et}}\).  The integer \(c\), the selected
site objects, and all these uniform vanishing assertions are part of the
datum.
\end{definition}

\begin{theorem}[Uniformly bounded-below proper coproducts]
Let \(p:Z\to Y\) be proper and let \((K_i)_{i\in I}\) be uniformly bounded
below.  The canonical comparison

\[
 \bigoplus_{i\in I}Rp_*K_i\longrightarrow
 Rp_*\!\left(\bigoplus_{i\in I}K_i\right)
\]

is an isomorphism.  Consequently the same assertion holds for
\(Rf_!=Rp_*j_!\) on every uniformly bounded-below family.
\end{theorem}

\begin{theorem}[Left completeness of the pro-etale derived category]
The standard t-structure on \(D(Z,\Lambda)\) is left complete: for every
\(K\), the canonical map

\[
 K\longrightarrow R\!\varprojlim_{N}\tau_{\geq-N}K
\]

is an isomorphism.  A map in \(D(Z,\Lambda)\) is therefore an isomorphism if
all its truncations \(\tau_{\geq-N}\) are isomorphisms.
\end{theorem}

\begin{theorem}[Full-dimension upgrade to arbitrary coproducts]
If \(\operatorname{FullCD}(p,\Lambda;c)\) holds, then for every set-indexed
family \((K_i)\) in \(D(Z,\Lambda)\), without a common lower bound, the map

\[
 \bigoplus_iRp_*K_i\xrightarrow{\sim}Rp_*\!\left(\bigoplus_iK_i\right)
\]

is an isomorphism.  This comparison is compatible with truncation.  If
\(Z\xrightarrow pY\xrightarrow qT\) is a composite, its compatibility with
the two-step proper-direct-image comparison is asserted only after separate
data \(\operatorname{FullCD}(p,\Lambda;c_p)\),
\(\operatorname{FullCD}(q,\Lambda;c_q)\), and
\(\operatorname{FullCD}(qp,\Lambda;c_{qp})\) have been supplied.
\end{theorem}

\begin{convention}[The explicit full-dimension hypothesis]
For every ambient representability statement below, choose a compactification
\(X\xrightarrow j\overline X\xrightarrow pY\), an integer \(c\), and assume
the statement \(\operatorname{FullCD}(p,\Lambda;c)\) of
Definition~\ref{def:full-cd}.
This is a hypothesis about all pro-etale \(\Lambda\)-module sheaves; a bound
restricted to constructible sheaves is a strictly weaker condition.
\end{convention}

\begin{corollary}[Coproducts for compact support and exceptional Hom]
Let \(X\xrightarrow j\overline X\xrightarrow pY\) be the chosen
compactification and assume \(\operatorname{FullCD}(p,\Lambda;c)\).  Then
\(Rf_!\) commutes with arbitrary coproducts.  Therefore, for every
\(L\in D(Y,\Lambda)\),

\[
 H_{f,L}\!\left(\bigoplus_{i\in I}K_i\right)
 \xrightarrow{\sim}\prod_{i\in I}H_{f,L}(K_i)
\]

is an isomorphism.
\end{corollary}

\begin{definition}[Brown relation-killing telescope]
Let \(H:\mathcal T^{\mathrm{op}}\to\mathsf{Ab}\) be cohomological and send
coproducts to products, where \(\mathcal T\) is generated by a shift-stable
set \(\mathcal G\) of compact objects.  Put

\[
 B_0=\bigoplus_{(G,\xi),\,G\in\mathcal G,\,\xi\in H(G)}G
\]

and choose \(e_0\in H(B_0)\) restricting to \(\xi\) on the corresponding
summand.  Given \((B_m,e_m)\), let

\[
 A_m=\bigoplus_{(G,u),\,H(u)(e_m)=0}G,
 \qquad a_m:A_m\longrightarrow B_m
\]

have component \(u\).  Complete \(a_m\) to a distinguished triangle
\(A_m\to B_m\xrightarrow{b_m}B_{m+1}\to A_m[1]\), and choose
\(e_{m+1}\in H(B_{m+1})\) with \(H(b_m)(e_{m+1})=e_m\).  If
\(s:\bigoplus_mB_m\to\bigoplus_mB_m\) has components
\(s\iota_m=\iota_{m+1}b_m\), define

\[
 B_H=\operatorname{Cone}\!\left(
 1-s:\bigoplus_mB_m\longrightarrow\bigoplus_mB_m\right).
\]
\end{definition}

\begin{theorem}[Brown telescope representability]
The compatible classes \(e_m\) induce \(e_H\in H(B_H)\).  The natural map

\[
 \operatorname{Hom}_{\mathcal T}(K,B_H)\longrightarrow H(K),
 \qquad u\longmapsto H(u)(e_H),
\]

is an isomorphism first for \(K\in\mathcal G\), then on the localizing
subcategory generated by \(\mathcal G\), and therefore for every
\(K\in\mathcal T\).  The pair \((B_H,e_H)\) is unique up to a unique
isomorphism compatible with the displayed transformation.
\end{theorem}

\begin{corollary}[Existence of the ambient finite-level right adjoint]
Assume \(\operatorname{FullCD}(p,\Lambda;c)\) for the chosen
compactification.  For every \(L\in D(Y,\Lambda)\), apply the Brown telescope
theorem to \(H_{f,L}\).
There is an object \(Rf^!L\in D(X,\Lambda)\) and a natural bijection

\[
 \Phi_{K,L}:
 \HomD_{D(Y,\Lambda)}(Rf_!K,L)
 \xrightarrow{\sim}
 \HomD_{D(X,\Lambda)}(K,Rf^!L).
\]

For a morphism \(v:L\to L'\), the unique map \(Rf^!(v):Rf^!L\to Rf^!L'\)
is the map representing postcomposition by \(v\).  These assignments define
an exact triangulated functor

\[
 Rf^!:D(Y,\Lambda)\longrightarrow D(X,\Lambda)
\]

right adjoint to \(Rf_!\).
\end{corollary}

\begin{definition}[Unit and counit]
The unit and counit are the morphisms

\[
 \eta_f:\operatorname{id}_{D(X,\Lambda)}\longrightarrow Rf^!Rf_!,
 \qquad
 \epsilon_f:Rf_!Rf^!\longrightarrow\operatorname{id}_{D(Y,\Lambda)}
\]

corresponding respectively to \(\operatorname{id}_{Rf_!K}\) and
\(\operatorname{id}_{Rf^!L}\) under \(\Phi\).  They satisfy

\[
 (\epsilon_fRf_!)\circ(Rf_!\eta_f)=\operatorname{id}_{Rf_!},
 \qquad
 (Rf^!\epsilon_f)\circ(\eta_fRf^!)=\operatorname{id}_{Rf^!}.
\]
\end{definition}

\begin{theorem}[Constructibility and amplitude of the Brown object]
Let \(f:X\to Y\) be separated of finite type, and assume
\(\operatorname{FullCD}(p,\Lambda;c)\) for the proper part of the chosen
compactification.  Assume also the constructible exceptional-devissage
theorem for \(f\), with finite cohomological bounds uniform on one finite
constructible stratification of \(X\).  Then the devissage datum supplies
nonnegative integers \(c_f^-\) and \(c_f^+\), depending only on \(f\), the
coefficient ring, and the selected stratifications, and independent of
\(a\), \(b\), and \(L\), such that

\[
 L\in D_c^{[a,b]}(Y,\Lambda)
 \quad\Longrightarrow\quad
 Rf^!L\in D_c^{[a-c_f^-,\,b+c_f^+]}(X,\Lambda).
\]

For a smooth morphism of pure relative dimension \(d\), pullback is
t-exact and purity gives the sharper interval

\[
 Rf^!D_c^{[a,b]}(Y,\Lambda)
 \subset D_c^{[a-2d,b-2d]}(X,\Lambda),
\]

so the displayed two-sided bound permits
\((c_f^-,c_f^+)=(2d,0)\).  For a regular closed immersion of codimension
\(r\), absolute purity gives, on complexes lisse and of finite Tor-amplitude
in a neighbourhood of the closed subscheme, the sharper interval

\[
 Ri^!D_c^{[a,b]}(Y,\Lambda)
 \subset D_c^{[a+2r,b+2r]}(X,\Lambda),
\]

so the two-sided bound permits \((c_i^-,c_i^+)=(0,2r)\).  No one-sided
amplitude assertion is made for a general separated finite-type morphism.
\end{theorem}

\begin{theorem}[Finite-Tor preservation and restricted adjunctions]
Under the hypotheses of the preceding theorem, suppose in addition that the
stratified finite-Tor devissage theorem holds for \(f\).  Then

\[
 L\in\Dctf(Y,\Lambda)\quad\Longrightarrow\quad
 Rf^!L\in\Dctf(X,\Lambda).
\]

Consequently the ambient adjunction restricts to exact adjunctions

\[
 Rf_!: \Dbc(X,\Lambda)\rightleftarrows\Dbc(Y,\Lambda):Rf^!,
\]

\[
 Rf_!: \Dctf(X,\Lambda)\rightleftarrows\Dctf(Y,\Lambda):Rf^!.
\]
\end{theorem}

\begin{convention}[Scope of the remaining exceptional formalism]
At finite coefficients, attach data separately to every proper map occurring
as the proper part of a chosen compactification, a common refinement, a
composite compactification, or a base-changed compactification.  For each
such map \(p':Z'\to Y'\), choose an integer \(c_{p'}\) and assume
\(\operatorname{FullCD}(p',\Lambda;c_{p'})\).  No FullCD datum is inferred
from the datum for another proper map.

Likewise, every occurrence on \(\Dbc\) assumes the constructible-amplitude
theorem separately for each compactifiable morphism in the diagram, and
every occurrence on \(\Dctf\) assumes its finite-Tor devissage theorem.
At adic or field coefficients, finite-reduction and scalar-extension
hypotheses are supplied for every coefficient stage and every functor in the
diagram.  The words ``composition'', ``base change'', and ``refinement''
below always include these separately quantified data; they do not assert
automatic propagation of any admissibility hypothesis.
\end{convention}

\section{Compactification formula and mates}

\begin{definition}[Mate of a natural transformation]
Let \(F_1,F_2:\mathcal C\to\mathcal D\) have right adjoints
\(G_1,G_2:\mathcal D\to\mathcal C\), with units \(\eta_i\) and counits
\(\epsilon_i\).  For a natural transformation
\(\theta:F_1\to F_2\), its right-adjoint mate is

\[
 \theta^\flat:G_2\longrightarrow G_1,
\]

defined by the typed composite

\[
 G_2\xrightarrow{\eta_1G_2}G_1F_1G_2
 \xrightarrow{G_1\theta G_2}G_1F_2G_2
 \xrightarrow{G_1\epsilon_2}G_1.
\]

Taking mates is bijective, reverses vertical composition, preserves
identities and pasting, and carries an isomorphism to an isomorphism.
\end{definition}

\begin{theorem}[Compactification formula]
Let \(c=(X\xrightarrow{j}\overline X\xrightarrow{p}Y)\) compactify \(f\).
Assume \(Rp_*:D(\overline X,\Lambda)\to D(Y,\Lambda)\) has a right adjoint

\[
 Rp^!:D(Y,\Lambda)\longrightarrow D(\overline X,\Lambda).
\]

Then \(j^*Rp^!\) is right adjoint to \(Rp_*j_!\), through the natural chain

\[
 \HomD_Y(Rp_*j_!K,L)
 \cong\HomD_{\overline X}(j_!K,Rp^!L)
 \cong\HomD_X(K,j^*Rp^!L).
\]

There is therefore a unique adjunction-compatible isomorphism

\[
 \kappa_c:Rf^!\xrightarrow{\sim}j^*Rp^!.
\]
\end{theorem}

\begin{theorem}[Independence of compactification]
For compactifications \(c,c'\), let
\(\rho_{c,c'}:Rf_{!,c}\xrightarrow{\sim}Rf_{!,c'}\) be the canonical
comparison obtained from a common refinement.  The inverse of its
right-adjoint mate is an isomorphism

\[
 \rho^!_{c,c'}:Rf^!_c\xrightarrow{\sim}Rf^!_{c'}.
\]

These maps satisfy

\[
 \rho^!_{c,c}=\operatorname{id},\qquad
 \rho^!_{c',c''}\circ\rho^!_{c,c'}=\rho^!_{c,c''},
\]

and commute with \(\eta_f\), \(\epsilon_f\), and the compactification
isomorphisms \(\kappa_c\).
\end{theorem}

\begin{corollary}[Proper and open normalizations]
If \(f\) is proper, then \(Rf_!=Rf_*\) and \(Rf^!\) is the right adjoint of
\(Rf_*\).  If \(j:U\hookrightarrow X\) is open, then

\[
 Rj^!\xrightarrow{\sim}j^*:D(X,\Lambda)\longrightarrow D(U,\Lambda),
\]

and the exceptional adjunction is \(j_!\dashv j^*\).
\end{corollary}

\begin{definition}[Exceptional base-change mate]
For a Cartesian square

\[
\begin{CD}
X' @>{u}>> X\\
@V{f'}VV @VV{f}V\\
Y' @>{v}>> Y,
\end{CD}
\]

compact-support base change gives

\[
 \operatorname{BC}_!:Lv^*Rf_!\longrightarrow Rf'_!Lu^*.
\]

Its right-adjoint mate is

\[
 Ru_*Rf'^!\longrightarrow Rf^!Rv_*.
\]

Whenever \(\operatorname{BC}_!\) is an isomorphism, the inverse mate is the
exceptional base-change isomorphism

\[
 \operatorname{BC}^!:Rf^!Rv_*\xrightarrow{\sim}Ru_*Rf'^!.
\]
\end{definition}

\begin{theorem}[Pasting of exceptional base change]
If compact-support base change is invertible for two composable Cartesian
squares, then exceptional base change is invertible for their composite.
The exceptional comparison for the pasted square equals the pasted
comparison for the two squares, after inserting the composition constraints
for direct images and exceptional inverse images.
\end{theorem}

\section{Composition, trace, and tensor exchange}

\begin{theorem}[Composition of exceptional inverse images]
Let \(X\xrightarrow fY\xrightarrow gZ\) be compactifiable.  The
right-adjoint mate of

\[
 \alpha_{g,f}:Rg_!Rf_!\xrightarrow{\sim}R(gf)_!
\]

is a canonical isomorphism

\[
 c_{g,f}:R(gf)^!\xrightarrow{\sim}Rf^!Rg^!.
\]

For three composable morphisms, the two maps from \(R(hgf)^!\) to
\(Rf^!Rg^!Rh^!\) coincide.  The identity comparison
\(R(\operatorname{id}_X)^!\cong\operatorname{id}\) satisfies both unit
laws.  All these maps are compatible with compactification comparison.
\end{theorem}

\begin{definition}[Exceptional trace]
The exceptional trace of \(f\) is the counit

\[
 \operatorname{Tr}_f=\epsilon_f:Rf_!Rf^!\longrightarrow
 \operatorname{id}_{D(Y,\Lambda)}.
\]
\end{definition}

\begin{proposition}[Trace transitivity]
For \(X\xrightarrow fY\xrightarrow gZ\), the equality

\[
 \operatorname{Tr}_{gf}\circ
 R(gf)_!(c_{g,f}^{-1})\circ\alpha_{g,f}
 =\operatorname{Tr}_g\circ Rg_!(\operatorname{Tr}_f)
\]

holds as a transformation \(Rg_!Rf_!Rf^!Rg^!\to\operatorname{id}\).
\end{proposition}

\begin{theorem}[Projection formula]
For \(K\in\Dctf(Y,\Lambda)\) and \(L\in\Dbc(X,\Lambda)\), the canonical map

\[
 Rf_!(Lf^*K\otimes^L L)
 \longrightarrow K\otimes^L Rf_!L
\]

is an isomorphism.  The projection isomorphism, the trace, and tensor--Hom
adjunction define
an exceptional exchange map

\[
 Rf^!R\mathcal Hom_Y(K,M)
 \longrightarrow
 R\mathcal Hom_X(Lf^*K,Rf^!M).
\]

This map is an isomorphism when \(K\) is perfect.
\end{theorem}

\section{Adic, rational, and algebraically closed coefficients}

\begin{theorem}[Finite reduction of compact support]
For every \(m\leq n\), coefficient extension and compact support carry a
canonical natural transformation on the ambient derived category,

\[
 \theta_{n,m,f}:\red_{n,m,Y}Rf_{!,\Lambda_n}
 \longrightarrow Rf_{!,\Lambda_m}\red_{n,m,X}.
\]

For every \(K\in\Dctf(X,\Lambda_n)\), its component
\(\theta_{n,m,f,K}\) is an isomorphism.  These transformations are natural
in \(K\), transitive in \(m\leq n\), and compatible with compactification
and composition.
\end{theorem}

\begin{definition}[Restriction and reduction mates]
On the ambient categories, the right-adjoint mate of \(\theta_{n,m,f}\)
with respect to
\(\red_{n,m}\dashv\operatorname{res}_{m,n}\) and
\(Rf_!\dashv Rf^!\) is the restriction mate

\[
 \theta^\flat_{n,m,f}:
 \operatorname{res}_{m,n,X}Rf^!_{\Lambda_m}
 \longrightarrow Rf^!_{\Lambda_n}\operatorname{res}_{m,n,Y}.
\]

For \(L\in\Dctf(Y,\Lambda_n)\), define the reduction mate

\[
 \beta_{n,m,f,L}:\red_{n,m,X}(Rf^!_{\Lambda_n}L)
 \longrightarrow
 Rf^!_{\Lambda_m}(\red_{n,m,Y}L).
\]

It is the unique morphism whose adjoint under
\(Rf_{!,\Lambda_m}\dashv Rf^!_{\Lambda_m}\) is

\[
\begin{aligned}
 Rf_{!,\Lambda_m}\red_{n,m,X}Rf^!_{\Lambda_n}L
 &\xrightarrow{\theta_{n,m,f,Rf^!L}^{-1}}
 \red_{n,m,Y}Rf_{!,\Lambda_n}Rf^!_{\Lambda_n}L\\
 &\xrightarrow{\red_{n,m,Y}(\epsilon_{f,L})}
 \red_{n,m,Y}L.
\end{aligned}
\]
Neither invertibility of the ambient restriction mate nor invertibility of
\(\beta_{n,m,f,L}\) is part of this definition.  The latter is the
finite-Tor coefficient-change statement below.
\end{definition}

\begin{theorem}[Proper exceptional coefficient change]
Let \(p:Z\to Y\) be proper.  Supply separate FullCD data, constructible
amplitude, and finite-Tor devissage at the \(\Lambda_n\)- and
\(\Lambda_m\)-stages.  Then, for every
\(L\in\Dctf(Y,\Lambda_n)\),

\[
 \beta_{n,m,p,L}:\red_{n,m,Z}Rp^!_{\Lambda_n}L
 \xrightarrow{\sim}
 Rp^!_{\Lambda_m}\red_{n,m,Y}L
\]

is an isomorphism.  It is natural in \(L\), transitive for
\(r\leq m\leq n\), and compatible with proper composition, units, counits,
and separately admissible proper base-change squares.
\end{theorem}

\begin{proposition}[Open exceptional coefficient change]
For a quasi-compact open immersion \(j:U\hookrightarrow X\), the
normalization \(Rj^!\cong j^*\) identifies
\(\beta_{n,m,j,L}\) with the canonical associativity isomorphism

\[
 \red_{n,m,U}j^*L\xrightarrow{\sim}
 j^*\red_{n,m,X}L.
\]

These isomorphisms are transitive in the coefficient indices and compatible
with composition of open immersions.
\end{proposition}

\begin{theorem}[Compactification devissage for the reduction mate]
Let \(f=pj\) be a chosen compactification.  Under the separately supplied
admissibility data for \(p\), the compactification formula
\(Rf^!\cong j^*Rp^!\) identifies \(\beta_{n,m,f,L}\) with the composite of
the proper comparison for \(p\) and the open comparison for \(j\).
Consequently

\[
 \beta_{n,m,f,L}:\red_{n,m,X}Rf^!_{\Lambda_n}L
 \xrightarrow{\sim}
 Rf^!_{\Lambda_m}\red_{n,m,Y}L
\]

is an isomorphism, independent of compactification.  These maps commute with
units, counits, separately admissible base-change mates, and the composition
maps \(c_{g,f}\); for \(r\leq m\leq n\), their two coefficient-reduction
composites agree under
\(\red_{m,r}\red_{n,m}\cong\red_{n,r}\).
\end{theorem}

\begin{definition}[Integral compact support in the system category]
For a compactification
\(c=(X\xrightarrow j\overline X\xrightarrow pY)\), extension by zero and
direct image act levelwise on the abelian system categories.  Derive \(p_*\)
by a functorial K-injective replacement in
\(\mathsf{Mod}(\overline X,\Lambda_\bullet)\), and set

\[
 Rf_{!,\OO_E,c}=Rp_*j_!:
 D(X,\Lambda_\bullet)\longrightarrow D(Y,\Lambda_\bullet).
\]

The system-level refinement maps define an exact triangulated functor
\(Rf_{!,\OO_E}\), independent of \(c\), together with
\(Rf_{!,\OO_E}\cong Rf_{!,\OO_E,c}\).
\end{definition}

\begin{theorem}[Finite levels of integral compact support]
Assume that the compact-support cohomological and finite-Tor bounds are
independent of \(n\).  The preceding functor restricts to

\[
 Rf_{!,\OO_E}:\Dctf(X,\OO_E)\longrightarrow\Dctf(Y,\OO_E),
\]

and for every \(K\in\Dctf(X,\OO_E)\) there are natural isomorphisms

\[
 \gamma_{n,f,K}:\red_{n,Y}Rf_{!,\OO_E}K
 \xrightarrow{\sim}
 Rf_{!,\Lambda_n}\red_{n,X}K.
\]

They commute with the normalization maps, compactification comparison,
composition, and finite reduction.
\end{theorem}

\begin{definition}[The common category of coefficient systems]
Let

\[
 \mathcal A_X(\Lambda_\bullet)
 =\mathsf{Mod}(X,\Lambda_\bullet)
\]

be the abelian category whose objects are families \(M_\bullet=(M_n)_{n\geq
1}\), with \(M_n\) a sheaf of \(\Lambda_{n,X}\)-modules, together with maps

\[
 t_{n,m}:M_n\longrightarrow
 \operatorname{res}_{m,n,X}M_m\qquad(m\leq n)
\]

that are linear through \(\Lambda_n\to\Lambda_m\), satisfy
\(t_{n,n}=\operatorname{id}\), and satisfy

\[
 t_{n,r}=\operatorname{res}_{m,n,X}(t_{m,r})\circ t_{n,m}
 \qquad(r\leq m\leq n)
\]

under
\(\operatorname{res}_{m,n,X}\operatorname{res}_{r,m,X}
=\operatorname{res}_{r,n,X}\).  Its ordinary derived category

\[
 D(X,\Lambda_\bullet)=D\bigl(\mathcal A_X(\Lambda_\bullet)\bigr)
\]

is the single coefficient category in which an integral exceptional object
will live.  Evaluation at level \(n\) is the exact triangulated functor
\(\operatorname{ev}_{n,X}:D(X,\Lambda_\bullet)\to D(X,\Lambda_n)\).
\end{definition}

\begin{definition}[Finite-stage exceptional tower]
Let \(L\in\Dctf(Y,\OO_E)\), put
\(L_n=\operatorname{ev}_{n,Y}L\), and write its normalization maps as

\[
 \nu^L_{n,m}:\red_{n,m,Y}L_n\xrightarrow{\sim}L_m
 \qquad(m\leq n).
\]

Put \(M_n=Rf^!_{\Lambda_n}L_n\).  The reduction transition is the typed
isomorphism

\[
\begin{aligned}
 \tau^L_{n,m}:\red_{n,m,X}M_n
 &\xrightarrow{\beta_{n,m,f,L_n}}
 Rf^!_{\Lambda_m}\red_{n,m,Y}L_n\\
 &\xrightarrow{Rf^!_{\Lambda_m}(\nu^L_{n,m})}M_m.
\end{aligned}
\]

Using \(\red_{n,m,X}\dashv\operatorname{res}_{m,n,X}\), let

\[
 \bar\tau^L_{n,m}:M_n\longrightarrow
 \operatorname{res}_{m,n,X}M_m
\]

be the adjoint of \(\tau^L_{n,m}\).  Thus \(\tau^L_{n,m}\) and
\(\bar\tau^L_{n,m}\) are morphisms in the explicitly displayed categories
\(D(X,\Lambda_m)\) and \(D(X,\Lambda_n)\), respectively.
\end{definition}

\begin{proposition}[Transition identities]
For \(r\leq m\leq n\), the reduction transitions satisfy

\[
 \tau^L_{n,r}=
 \tau^L_{m,r}\circ\red_{m,r,X}(\tau^L_{n,m})
\]

after the canonical reassociation
\(\red_{m,r,X}\red_{n,m,X}\cong\red_{n,r,X}\), and
\(\tau^L_{n,n}=\operatorname{id}\).  Equivalently, their restriction
adjoints satisfy

\[
 \bar\tau^L_{n,r}=
 \operatorname{res}_{m,n,X}(\bar\tau^L_{m,r})
 \circ\bar\tau^L_{n,m},
 \qquad \bar\tau^L_{n,n}=\operatorname{id}.
\]

For every \(v:L\to L'\), the maps
\(Rf^!_{\Lambda_n}(v_n):M_n\to M'_n\) commute with both forms of the
transition maps.  These equalities remain valid after inserting the
finite-level units, counits, compactification comparisons, or composition
maps.
\end{proposition}

\begin{definition}[Uniform integral exceptional-amplitude datum]
Such a datum for \(f\) consists of nonnegative integers
\(c_f^-,c_f^+,t_f^-,t_f^+\) and, for every finite constructible
stratification \(\mathcal S\) of \(Y\), a finite constructible
stratification \(f^!\mathcal S\) of \(X\), all independent of \(n\), with
the following property.  If \(L_n\in\Dctf(Y,\Lambda_n)\) is constructible
along \(\mathcal S\), lies in \(D_c^{[a,b]}\), and has Tor-amplitude in
\([r,s]\), then

\[
 Rf^!_{\Lambda_n}L_n\in
 D_c^{[a-c_f^-,\,b+c_f^+]}(X,\Lambda_n)
\]

is constructible along \(f^!\mathcal S\) and has Tor-amplitude in
\([r-t_f^-,s+t_f^+]\).
\end{definition}

\begin{definition}[Selected realization category]
Fix a uniform integral exceptional-amplitude datum.  Define
\(\mathsf{ExcSys}_f(Y,\OO_E)\) as follows.  An object is a triple
\((L,M,\sigma)\), where \(L\in\Dctf(Y,\OO_E)\),
\(M\in D(X,\Lambda_\bullet)\), and

\[
 \sigma_{n}:\operatorname{ev}_{n,X}M\xrightarrow{\sim}M_n
 =Rf^!_{\Lambda_n}L_n
\]

is an isomorphism for every \(n\).  It is required that, for \(m\leq n\),

\[
\begin{CD}
\red_{n,m,X}\operatorname{ev}_{n,X}M
 @>{\red_{n,m,X}(\sigma_n)}>>\red_{n,m,X}M_n\\
@V{\nu^M_{n,m}}V{\sim}V @VV{\tau^L_{n,m}}V\\
\operatorname{ev}_{m,X}M @>{\sigma_m}>>M_m
\end{CD}
\]

commute.  A morphism
\((v,w):(L,M,\sigma)\to(L',M',\sigma')\) consists of
\(v:L\to L'\) and \(w:M\to M'\), with

\[
 \sigma'_n\circ\operatorname{ev}_{n,X}(w)
 =Rf^!_{\Lambda_n}(v_n)\circ\sigma_n
\]

for every \(n\).  The forgetful functor

\[
 \pi_f:\mathsf{ExcSys}_f(Y,\OO_E)\longrightarrow\Dctf(Y,\OO_E),
 \qquad(L,M,\sigma)\longmapsto L,
\]

is exact for the componentwise distinguished triangles.
\end{definition}

\begin{theorem}[Realization of the finite-stage tower]
The ordinary derived category of the abelian system category admits an exact
triangulated section

\[
 s_f:\Dctf(Y,\OO_E)\longrightarrow\mathsf{ExcSys}_f(Y,\OO_E),
 \qquad \pi_fs_f=\operatorname{id}.
\]

If \(s_f(L)=(L,M,\sigma)\), then \(M\) is normalized, is bounded
constructible with the common interval
\([a-c_f^-,b+c_f^+]\), and has the common Tor interval
\([r-t_f^-,s+t_f^+]\) whenever \(L\) has bounds \([a,b]\) and \([r,s]\).
The selected sections for all morphisms may be chosen so that the finite
stage unit, counit, compactification, and composition rectangles are equal
after every evaluation functor.  This is a statement about ordinary
triangulated categories and the derived categories of the displayed abelian
system categories.
\end{theorem}

\begin{definition}[Integral exceptional inverse image]
Choose a section \(s_f\), write
\(s_f(L)=(L,M_f(L),\sigma_L)\), and define

\[
 Rf^!_{\OO_E}(L)=M_f(L),\qquad
 Rf^!_{\OO_E}(v)=\text{the second component of }s_f(v).
\]

Thus

\[
 Rf^!_{\OO_E}:\Dctf(Y,\OO_E)\longrightarrow\Dctf(X,\OO_E)
\]

is an exact triangulated functor, and its finite-level comparison is
\(\sigma_{n,L}:\operatorname{ev}_{n,X}Rf^!_{\OO_E}L
\xrightarrow{\sim}Rf^!_{\Lambda_n}L_n\).
\end{definition}

\begin{definition}[Derived Hom and derived inverse limit]
For \(C=\Lambda_n\), let
\(\mathcal A_X(C)=\mathsf{Mod}(X_{\mathrm{pro\acute et}},C_X)\), and let
\(\mathcal A_X(\Lambda_\bullet)=\mathsf{Mod}(X,\Lambda_\bullet)\).
Define

\[
 \mathbf R\!\operatorname{Hom}_{X,C}:
 D(\mathcal A_X(C))^{\mathrm{op}}\times D(\mathcal A_X(C))
 \longrightarrow D(\mathsf{Ab})
\]

as the total right-derived functor of the ordinary Hom bifunctor.  It is
computed by a functorial K-injective replacement in the second variable and
satisfies

\[
 H^q\mathbf R\!\operatorname{Hom}_{X,C}(K,M)
 \cong\operatorname{Hom}_{D(\mathcal A_X(C))}(K,M[q]).
\]

For a tower \((C_n,t_{n+1,n})\) of cochain complexes of abelian groups,
let \(R\!\varprojlim_n C_n\in D(\mathsf{Ab})\) be the ordinary right-derived
inverse limit in the abelian category of such towers.  For a countable tower
it is characterized by the distinguished triangle

\[
 R\!\varprojlim_n C_n\longrightarrow
 \prod_{n\geq1}C_n\xrightarrow{1-s}\prod_{n\geq1}C_n
 \longrightarrow (R\!\varprojlim_n C_n)[1],
\]

where the \(n\)-th component of \(s\) is
\(t_{n+1,n}:C_{n+1}\to C_n\).
\end{definition}

\begin{theorem}[Adic derived-Hom comparison]
Let \(K,M\in\Dctf(X,\OO_E)\).  The functorial K-injective resolutions,
finite reduction, and the normalization maps define a tower of derived-Hom
complexes representing

\[
 \bigl(\mathbf R\!\operatorname{Hom}_{X,\Lambda_n}(K_n,M_n)\bigr)_{n\geq1}
\]

in \(D(\mathsf{Ab})\).  There is a natural isomorphism

\[
 \mathbf R\!\operatorname{Hom}_{X,\Lambda_\bullet}(K,M)
 \xrightarrow{\sim}
 R\!\varprojlim_n
 \mathbf R\!\operatorname{Hom}_{X,\Lambda_n}(K_n,M_n).
\]

For every \(q\in\mathbf Z\), its cohomology gives the natural exact sequence

\[
\begin{aligned}
0\longrightarrow{}&\varprojlim
{}^{1}_n\operatorname{Hom}_{D(X,\Lambda_n)}(K_n,M_n[q-1])\\
\longrightarrow{}&\operatorname{Hom}_{D(X,\Lambda_\bullet)}(K,M[q])\\
\longrightarrow{}&\varprojlim_n
\operatorname{Hom}_{D(X,\Lambda_n)}(K_n,M_n[q])
\longrightarrow0.
\end{aligned}
\]
\end{theorem}

\begin{theorem}[Adic adjunction and reduction rectangles]
Let \(K\in\Dctf(X,\OO_E)\) and \(L\in\Dctf(Y,\OO_E)\).  For every
\(q\in\mathbf Z\), the finite-level adjunction isomorphisms

\[
 \operatorname{Hom}_{D(Y,\Lambda_n)}
 (Rf_{!,\Lambda_n}K_n,L_n[q])
 \xrightarrow{\sim}
 \operatorname{Hom}_{D(X,\Lambda_n)}
 (K_n,Rf^!_{\Lambda_n}L_n[q])
\]

form an isomorphism of towers under \(\gamma_{n,f,K}\),
\(\sigma_{n,L}\), and the finite coefficient-change rectangles.  They
identify the \(\varprojlim\)- and \(\varprojlim^1\)-terms in the two adic
derived-Hom exact sequences.  The middle comparison is a natural isomorphism

\[
 \HomD_{\Dctf(Y,\OO_E)}(Rf_{!,\OO_E}K,L)
 \xrightarrow{\sim}
 \HomD_{\Dctf(X,\OO_E)}(K,Rf^!_{\OO_E}L).
\]

Thus \(Rf_{!,\OO_E}\dashv Rf^!_{\OO_E}\).  For a morphism
\(v:L\to L'\) and every \(n\), the square

\[
\begin{CD}
\red_{n,X}Rf^!_{\OO_E}L @>{\sigma_{n,L}}>>
 Rf^!_{\Lambda_n}\red_{n,Y}L\\
@V{\red_{n,X}(Rf^!_{\OO_E}v)}VV
 @VV{Rf^!_{\Lambda_n}(\red_{n,Y}v)}V\\
\red_{n,X}Rf^!_{\OO_E}L' @>{\sigma_{n,L'}}>>
 Rf^!_{\Lambda_n}\red_{n,Y}L'
\end{CD}
\]

commutes.  The corresponding rectangles commute for units, counits,
composition, and compactification comparison.
\end{theorem}

\begin{theorem}[Scalar extension]
Let \(A\to B\) be \(\OO_E\to E\), a finite extension \(E\to E'\), or a
finite coefficient extension.  For \(L\in\Dctf(Y,A)\), the canonical mate
is an isomorphism

\[
 L\operatorname{Ext}_{B/A,X}(Rf^!_A L)
 \xrightarrow{\sim}
 Rf^!_B\bigl(L\operatorname{Ext}_{B/A,Y}L\bigr).
\]

It is transitive in towers \(A\to B\to C\), independent of an integral
lattice after rationalization, and compatible with all adjunction and
composition constraints.
\end{theorem}

\begin{definition}[The \(\Qlb\)-linear exceptional inverse image]
Represent an object of \(\Dctf(Y,\Qlb)\) by an idempotent
\((E,L,e)\).  Choose a finite extension \(F/E\) over which all data are
defined and set

\[
 Rf^!_{\Qlb}(E,L,e)=
 \bigl(F,Rf^!_F(L\otimes_EF),Rf^!_F(e\otimes_EF)\bigr).
\]

The scalar-extension theorem makes this independent of \(F\).  It defines

\[
 Rf^!_{\Qlb}:\Dctf(Y,\Qlb)\longrightarrow\Dctf(X,\Qlb).
\]
\end{definition}

\begin{theorem}[Coefficient-complete exceptional formalism]
Let the entire diagram under consideration carry the separately quantified
FullCD, constructible-amplitude, finite-Tor, finite-reduction, and
scalar-extension data of the scope convention.  At every finite, adic,
rational, and \(\Qlb\) coefficient stage covered by those data, the functor
\(Rf^!\) is
right adjoint to \(Rf_!\) on the displayed \(\Dctf\) categories.  Proper and
open normalization, compactification independence,
base-change mates, composition, trace transitivity, and tensor exchange are
preserved by reduction, rationalization, finite coefficient-field extension, and passage
to idempotent summands.
\end{theorem}

\end{document}
